\subsection{同伦方法}\label{sec:10e}
Floer 理论中还有另一种证明良定义性定理的方法。这种方法相当直接，历史上也早于延拓映射，但在具体实施时，技术上可能颇为棘手。虽然它并非我们当前目的所必需，我们仍介绍这一方法，因为它本身就很重要。

设对于某个带角正合辛流形 $M$，给定 Fukaya 范畴的两个版本 $\F(M)^{\mathrm{pr},s}$，$s=0,1$。假定两者采用相同的条带状端万有选择，因而区别来自所使用的 Floer 数据及扰动数据。进一步假设，对每一对 $(L_0,L_1)$，都存在 Floer 数据的光滑同伦
\begin{equation}\label{eq:10.10}
(H^s_{L_0,L_1},J^s_{L_0,L_1}),\qquad 0\leq s\leq1,
\end{equation}
连接定义 $\F(M)^{\mathrm{pr},0}$ 和 $\F(M)^{\mathrm{pr},1}$ 时使用的数据。注意，第一个假设只是技术性的，第二个假设却相当强，因为它排除了生灭现象。\znote{142}{生灭现象指参数 $s$ 经过某个临界值时，$\phi^{s,1}(L_0)$ 与 $L_1$ 由横截变为不横截（或反之），交点因此成对地产生或湮灭。若允许这种情形，交点集便不能随 $s$ 连续地等同，\eref{10.11}那样的同构也不存在，构造必须改用带生灭的参数族；而\eref{10.10}要求对每个 $s$ 都是 Floer 数据（交点横截），所以这里是较简单的情形。[来源：原书 §(10e) 本段；另一种处理见 §(8k) 的延拓映射。]}为解释这一点，固定 $(L_0,L_1)$，以下略去这些下标。设 $X^s$ 是对应于 $H^s$ 的 Hamilton 向量场，它还依赖另一个时间变量 $t\in[0,1]$。以 $\phi^{s,t}$ 表示沿后一个变量积分得到的 Hamilton 微分同胚族，即 $\phi^{s,0}=\Id_M$，且 $\partial\phi^{s,t}/\partial t=X^{s,t}(\phi^{s,t})$。由于所有\eref{10.10}都必须是 Floer 数据，交集 $\phi^{s,1}(L_0)\cap L_1$ 都是横截的，故交点随 $s$ 光滑变化。换言之，设 $CF^{\mathrm{pr}}(L_0,L_1)^s$ 是对某个 $s\in[0,1]$ 用\eref{10.10}定义的 Floer 上链群，$\mathcal C(L_0,L_1)^s$ 是其标准生成元集合，由满足 $y^s(0)\in L_0$、$y^s(1)\in L_1$ 及 $dy^s/dt=X^{s,t}(y^s(t))$ 的映射 $y^s:[0,1]\to M$ 组成。同样，以 $\mathcal C(L_0,L_1)^{[0,1]}$ 表示光滑映射 $y:[0,1]^2\to M$ 的集合，要求每个限制 $y^s=y(s,\cdot)$ 都属于 $\mathcal C(L_0,L_1)^s$；以 $CF^{\mathrm{pr}}(L_0,L_1)^{[0,1]}$ 表示该集合生成的 $\K$ 向量空间。于是，对任意 $s$，限制映射 $\mathcal C(L_0,L_1)^{[0,1]}\to \mathcal C(L_0,L_1)^s$ 都是双射，从而诱导同构
\begin{equation}\label{eq:10.11}
CF^{\mathrm{pr}}(L_0,L_1)^{[0,1]}\cong CF^{\mathrm{pr}}(L_0,L_1)^s.
\end{equation}
若以为仅凭这些同构便能得到 $\Ai$ 函子，就过于天真了。实际上，必须考虑加入 $s$ 参数的非齐次伪全纯映射的适当模空间。
\SourcePage{151}
先从一阶构造，也就是上链映射的构造开始。这个构造相当清楚，可以作为一般情形的模型。给定 $y_0,y_1\in \mathcal C(L_0,L_1)^{[0,1]}$，记 $\M_Z^*(y_0,y_1)^s=\M_Z^*(y_0^s,y_1^s)$，表示关于数据 $(H^s,J^s)$ 的 Floer 方程的非驻定解模空间，其极限为 $y_0^s,y_1^s$，并照例除去平移作用。同样，记 $\M_Z^*(y_0,y_1)^{[0,1]}$ 为二元组 $(s,u)$ 的空间，其中 $s\in[0,1]$，$u\in\M_Z^*(y_0,y_1)^s$。对于每个 $\nu\geq1$，现在引入相关的模空间 $\mathcal H^\nu(y_0,y_1)^{[0,1]}$。它的点是 $\nu$ 元组 $((s_1,u_1),\ldots,(s_\nu,u_\nu))$，满足 $s_1\leq\cdots\leq s_\nu$，且对某些 $y_0=\widetilde y_0,\widetilde y_1,\ldots,\widetilde y_\nu=y_1$，每个 $u_i\in\M_Z^*(\widetilde y_{i-1},\widetilde y_i)^{s_i}$。可以把 $\mathcal H^\nu(y_0,y_1)$ 的元素抽象地画成由双标点圆盘组成的链，每个圆盘带有一个变量 $s_k$，见\sr{fig:10.2}{图~10.2}。为书写方便，形式上规定：当 $y_0=y_1$ 时，$\mathcal H^0(y_0,y_1)^{[0,1]}$ 是单点；否则为空集。
\SourceFigure{figure-10-2}{0.58569168}{10.2}
可以将这一构造看作平凡族 $[0,1]\times Z\to[0,1]$ 的带参数模问题，只是与\secref{9h}的设置不同，端所对应的 Floer 数据也随参数 $s$ 变化。不过，基本结构仍然相同；特别地，每个 $(s,u)\in\M_Z^*(y_0,y_1)^{[0,1]}$ 都有对应的扩展线性化算子，它像\eref{9.18}那样将额外变量纳入考虑。正则性定义为这个算子的满射性，并蕴含模空间的局部光滑性。如果同伦\eref{10.10}选得足够一般，正则性便处处成立，以下假定如此。为了说明空间 $\mathcal H^\nu(y_0,y_1)^{[0,1]}$ 具有良好性质，还必须证明它们的边界面本身也正则。这些边界面是乘积
\[
\prod_{i=1}^\nu\M_Z^*(\widetilde y_{i-1},\widetilde y_i)^{[0,1]}
\]
中由一个或多个形如 $s_i=s_{i+1}$ 的方程截出的子流形。使用稍复杂一些的标准横截性论证，可证明只要各 $(s_i,u_i)$ 两两不同，就能实现这一点。幸运的是，在目前情形中这个条件自动满足，因为每个 $u_i$ 都使 $H^{s_i}$ 扰动后的作用量泛函严格下降。

能量论证表明，当 $\nu$ 充分大时，$\mathcal H^\nu(y_0,y_1)^{[0,1]}$ 为空。此外，对固定的 $\nu$，带额外 $s$ 参数的 Gromov 紧性说明，$\mathcal H^\nu(y_0,y_1)^{[0,1]}$ 中只有有限多个孤立点。因此可以
\SourcePage{152}
定义映射 $\Gamma:CF^{\mathrm{pr}}(L_0,L_1)^1\to CF^{\mathrm{pr}}(L_0,L_1)^0$，令
\begin{equation}\label{eq:10.12}
\Gamma(y_1^1)=\sum_{y_0}\sum_{\nu\geq0}\#\mathcal H^\nu(y_0,y_1)^{[0,1]}\,y_0^0.
\end{equation}
\znote{144}{这与\eref{9.23}的计数方式相同，只是多了一个参数 $s$：$\nu$ 元组 $(s_1,u_1),\ldots,(s_\nu,u_\nu)$ 相当于把一条链上各段 Floer 条带按 $s$ 排好，$\#$ 仍是零维模空间的模 $2$ 点数。$\nu$ 充分大时空间为空（能量论证），故外层对 $\nu$ 的求和有限；对每个 $\nu$，带 $s$ 参数的 Gromov 紧性又保证孤立点只有有限多个。这两点合起来使\eref{10.12}是良定义的 $\K$-线性映射。[来源：原书 §(10e) 与 §(9j) 的\eref{9.23}。]}
\begin{lemma}\label{st:10.12}
$\Gamma$ 是（不分次）链复形的同构。
\end{lemma}
\begin{proof}
证明思路很直接。记 $\mathcal H^\nu(y_0,y_1)^{[0,1],1}$ 为 $\mathcal H^\nu(y_0,y_1)^{[0,1]}$ 的所有一维连通分支的并。这是带边界的一维流形，边界点出现于以下情形：(i) $s_1=0$；(ii) $s_\nu=1$；或 (iii) 对某个 $i$，$s_{i-1}=s_i$。取 Gromov 紧化 $\overline{\mathcal H}^\nu(y_0,y_1)^{[0,1],1}$ 后，又加入一些边界点，对应于某个 $u_i$ 断裂为一对 Floer 轨线，即对某个 $\widetilde y$，得到 $\M_Z^*(\widetilde y_{i-1},\widetilde y)^{s_i}\times\M_Z^*(\widetilde y,\widetilde y_i)^{s_i}$ 的孤立点。在边界点的代数计数中，这些点恰好抵消上述情形 (iii) 的贡献，剩下等式
\[
\begin{aligned}
&\sum_{\nu\geq0}\sum_{\widetilde y}\#\M_Z^*(y_0,\widetilde y)^0\,\#\mathcal H^\nu(\widetilde y,y_1)^{[0,1]}\\
&\qquad+\sum_{\nu\geq0}\sum_{\widetilde y}\#\mathcal H^\nu(y_0,\widetilde y)^{[0,1]}\,\#\M_Z^*(\widetilde y,y_1)^1=0,
\end{aligned}
\]
这恰好说明 $\Gamma$ 是链映射。要证明它是同构，论证大致如下。如果模空间 $\M_Z^*(y_0,y_1)^{[0,1]}$ 没有孤立点，则由定义 $\Gamma=\Id$。再看次简单的情形，即孤立点恰好只出现在参数 $s$ 的一个取值处。这时 $\Gamma=\Id+(\text{下三角部分})$，其中第二项关于 $A_{H^s}$ 作用量的值所给出的滤过严格下降。最后，利用\eref{10.12}与同伦复合的良好相容性，将一般情形归结为已经讨论的情形。\znote{144}{“$\Gamma=\Id+$下三角”之所以给出同构，是因为按 $A_{H^s}$ 的作用量值把生成元排成有限降链后，$\Gamma$ 在该滤过下的矩阵对角元为 $1$，非对角元只出现在作用量严格更小处；这样的矩阵可逆，其逆按作用量差逐级（Neumann 级数）给出，每一级只有有限多项。这正是证明要从“无孤立点故 $\Gamma=\Id$”逐步放宽到一般情形的理由。[来源：原书 §(10e)；作用量滤过与\eref{8.12} 的能量恒等式相容。]}沿此思路的完整论证，见下文\stref{笔记}{10.14}所列参考文献。
\end{proof}
对于伪全纯 $(d+1)$ 边形（$d\geq2$），相应构造在技术上稍复杂一些。对每个正合 Lagrangian 子流形的 $(d+1)$ 元组 $(L_0,\ldots,L_d)$，选择
\begin{equation}\label{eq:10.13}
\sigma_{L_0,\ldots,L_d,k},\quad1\leq k\leq d,\qquad
(K^s_{L_0,\ldots,L_d},J^s_{L_0,\ldots,L_d}),\quad0\leq s\leq1,
\end{equation}
满足下述性质；讨论这些性质时，略去下标 $L_0,\ldots,L_d$。每个 $\sigma_k$ 是光滑函数 $\mathcal R^{d+1}\times[0,1]\to[0,1]$，当 $s=0,1$ 时满足 $\sigma_k(r,s)=s$，称为\emph{延迟函数}。$(K^s,J^s)$ 是万有族 $\mathcal S^{d+1}$ 上扰动数据的同伦，在端点 $s=0,1$ 应还原为 $\F(M)^{\mathrm{pr},s}$ 使用的数据，并且在由延迟函数确定的 $s$ 变量变换之后，与 Floer 数据\eref{10.10}相容。具体地，在每个条带状端 $\eps_k^{d+1}:\mathcal R^{d+1}\times Z^\pm\to\mathcal S^{d+1}$ 的像上，要求
\begin{equation}\label{eq:10.14}
(\eps_k^{d+1})^*(K^s,J^s)=
\begin{cases}
(H^s_{L_0,L_d}(t)\,dt,J^s_{L_0,L_d}(t)),&k=0,\\
(H^{\widetilde s}_{L_{k-1},L_k}(t)\,dt,J^{\widetilde s}_{L_{k-1},L_k}(t)),
&k>0,\quad\widetilde s=\sigma_k(r,s).
\end{cases}
\end{equation}
\SourcePage{153}
这里的记号有些冲突：$t$ 是半无穷条带 $Z^\pm$ 的第二个坐标，然而 $s$ 并不是第一个坐标（该坐标在下面的论证中不起作用），而是控制同伦的变量。入向端与出向端的 Floer 数据\eref{10.14}一般对应于不同的取值 $\widetilde s\ne s$，这一事实影响曲面的粘合方式。现在说明怎样将\secref{9i}最初引入的扰动数据一致性选择概念调整到这里。

设 $T$ 为稳定的 $d$ 叶树。对每个顶点 $v\in\operatorname{Ve}(T)$，固定一个 $|v|$ 标点圆盘 $S_v$，以 $r_v\in\mathcal R^{|v|}$ 表示对应的模参数。为这些圆盘配备彼此相容的 Lagrangian 标记。再固定 $s\in[0,1]$。随后可用下面的递归规则为每个顶点赋予参数值 $s_v\in[0,1]$：首先，若 $v$ 是离根最近的顶点，令 $s_v=s$。其次，设 $e$ 为连接顶点 $v^-(e)$ 与 $v^+(e)$ 的任意内边；回忆\secref{9d}，边的定向指向远离根的方向。假设对应的旗标 $f^+(e)=(v^-(e),e)$，按典范次序是邻接 $v^-(e)$ 的第 $k_e$ 个正旗标，则令
\begin{equation}\label{eq:10.15}
s_{v^+(e)}=\sigma_{k_e}(r_{v^-(e)},s_{v^-(e)}),
\end{equation}
其中 $\sigma_{k_e}$ 是以 $S_{v^-(e)}$ 的 $|v^-(e)|$ 个 Lagrangian 标记为指标的延迟函数之一。现在以粘合参数 $\{\rho_e\}$ 将各 $S_v$ 粘合起来，以 $r\in\mathcal R^{d+1}$ 表示所得 $(d+1)$ 标点圆盘 $S$ 的分类模点。回忆，$S$ 的出向无穷远点可与 $T$ 的叶子识别。若叶子对应的旗标为 $f_{k_1}(v_1),\ldots,f_{k_d}(v_d)$，就在这些无穷远点上分别赋予数值 $s^1=\sigma_{k_1}(r_{v_1},s_{v_1}),\ldots,s^d=\sigma_{k_d}(r_{v_d},s_{v_d})\in[0,1]$。一致性条件的第一部分只涉及延迟函数：当粘合参数充分小时，这些数值应分别等于 $S$ 自然对应的数值 $\sigma_1(r,s),\ldots,\sigma_d(r,s)$。等式两边都略去了 Lagrangian 标记，但选择都是自然的：由于已为所有顶点选定相容标记，每个 $S_{v_k}$ 都带有 Lagrangian 标记，而 $S$ 通过粘合继承这些标记。接着，在每个 $S_v$ 上配备时间 $s_v$ 处的扰动数据\eref{10.13}。条件\eref{10.14}与\eref{10.15}共同保证这些数据在对应条带状端上具有相同的行为，因而可以粘合，得到 $S$ 上的扰动数据。另一方面，也可以通过分类同构 $S\cong\mathcal S_r^{d+1}$ 拉回\eref{10.13}，为同一个圆盘配备扰动数据。一致性条件要求这两种数据满足与\secref{9i}相同的关系：当粘合参数充分小时，在 $S$ 的薄部完全相等，在厚部渐近一致；后一个条件用万有族的部分紧化来定义。

选定满足这种一致性条件的\eref{10.13}后，就能引入对应的模空间。给定 $(L_0,\ldots,L_d)$ 及 $y_0\in \mathcal C(L_0,L_d)^{[0,1]}$、$y_k\in \mathcal C(L_{k-1},L_k)^{[0,1]}$，令 $\M^{d+1}(y_0,\ldots,y_d)^{[0,1]}$ 为
\SourcePage{154}
三元组 $(r,s,u)$ 的空间，其中 $r\in\mathcal R^{d+1}$、$s\in[0,1]$，而 $u:\mathcal S_r^{d+1}\to M$ 满足采用扰动数据\eref{10.13}的非齐次伪全纯映射方程；其在入向条带状端上的极限为 $y_0^s$，在各出向端上的极限为 $y_k^{\widetilde s}$，其中 $\widetilde s=\sigma_{L_0,\ldots,L_d,k}(r,s)$。与前面考虑的 $\mathcal H^\nu(y_0,y_1)^{[0,1]}$ 类似，还存在另一族密切相关的模空间 $\mathcal H^T(y_0,\ldots,y_d)^{[0,1]}$。其定义部分地类似于\secref{9l}讨论的紧化，我们借用那里的某些记号。固定半稳定 $d$ 叶树 $T$。$\mathcal H^T(y_0,\ldots,y_d)^{[0,1]}$ 的一个点由下列数据组成：
\begin{itemize}
\item 对每个满足 $|v|\geq3$ 的顶点，给定 $r_v\in\mathcal R^{|v|}$，从而给定一个 $|v|$ 标点圆盘 $S_v$。对其余顶点，令 $S_v=Z$。
\end{itemize}
可以用唯一的方式给各 $S_v$ 配备相容的 Lagrangian 标记，使粘合后的 $(d+1)$ 标点圆盘上的诱导标记等于给定的 $L_k$。完成这一选择后，下一项数据是：
\begin{itemize}
\item 对每个 $v$，给定数值 $s_v\in[0,1]$，满足\eref{10.15}的不等式版本。也就是说，对于每条内边 $e$，
\begin{equation}\label{eq:10.16}
s_{v^+(e)}\geq
\begin{cases}
\sigma_{k_e}(r_{v^-(e)},s_{v^-(e)}),&|v^-(e)|\geq3,\\
s_{v^-(e)},&|v^-(e)|=2.
\end{cases}
\end{equation}
\end{itemize}
接着，回忆在这里每个旗标 $f$ 都对应一对 Lagrangian 子流形 $(L_{f,0},L_{f,1})$。
\begin{itemize}
\item 对每个 $f$，指定 $\widetilde y_f\in \mathcal C(L_{f,0},L_{f,1})^{[0,1]}$，满足下列条件。若 $f^-$ 是对应于 $T$ 的根的唯一旗标，则 $\widetilde y_{f^-}=y_0$。同样，对应于叶子的 $d$ 个旗标上的 $\widetilde y_f$ 分别为 $y_1,\ldots,y_d$。若 $f^\pm(e)$ 是一条内边的两个旗标，则 $\widetilde y_{f^-(e)}=\widetilde y_{f^+(e)}$。
\item 最后，对每个 $v$，给定映射 $u_v:S_v\to M$；当 $|v|=2$ 时，它是\eref{8.8}的非驻定解，当 $|v|\geq3$ 时，它是\eref{8.9}的解。相应的 Floer 数据和扰动数据分别由\eref{10.10}和\eref{10.13}在 $s=s_v$ 处给出；条带状端上的极限应为 $v$ 周围各旗标 $f$ 对应的 $\widetilde y_f$，在上述相应参数处取值。等价地，
\[
\begin{cases}
(s_v,u_v)\in\M_Z^*(\widetilde y_{f_0(v)},\widetilde y_{f_1(v)})^{[0,1]},&|v|=2,\\
(r_v,s_v,u_v)\in\M^{|v|}(\widetilde y_{f_0(v)},\ldots,\widetilde y_{f_{|v|-1}(v)})^{[0,1]},&|v|\geq3.
\end{cases}
\]
\end{itemize}
\sr{fig:10.3}{图~10.3}示意各曲面 $S_v$ 及其 Lagrangian 标记，并画出各 $s_v$ 之间相应的不等式。照例，若两组映射 $\{u_v:S_v\to M\}$ 只相差各不稳定分支 $S_v\cong Z$ 上的平移，则视为相同。
\SourcePage{155}
\SourceFigure{figure-10-3}{0.64592238}{10.3}
假定\eref{10.13}中的选择足够一般，就可以仿照\eref{10.12}定义多线性映射
\begin{equation}\label{eq:10.17}
\begin{aligned}
&\G^d:CF^{\mathrm{pr}}(L_{d-1},L_d)^1\otimes\cdots\otimes CF^{\mathrm{pr}}(L_0,L_1)^1
\longrightarrow CF^{\mathrm{pr}}(L_0,L_d)^0,\\
&\G^d(y_d^1,\ldots,y_1^1)=\sum_{y_0}\sum_T\#\mathcal H^T(y_0,\ldots,y_d)^{[0,1]}\,y_0^0.
\end{aligned}
\end{equation}
\begin{proposition}\label{st:10.13}
\eref{10.12}中的映射 $\G^1=\Gamma$，连同\eref{10.17}中的高阶项 $\G^d$，构成（不分次）$\Ai$ 函子
\[
\G:\F(M)^{\mathrm{pr},1}\longrightarrow\F(M)^{\mathrm{pr},0}.
\]
\end{proposition}
注意，由\stref{引理}{10.12}，线性部分 $\G^1$ 都是同构；若通过\eref{10.11}识别 $s=0,1$ 时的 Floer 上链空间 $CF^{\mathrm{pr}}(L_0,L_1)^s$，则 $\G$ 是\secref{1c}意义下的形式
\SourcePage{156}
微分同胚。特别地，$\F(M)^{\mathrm{pr},0}$ 与 $\F(M)^{\mathrm{pr},1}$ 是拟同构的 $\Ai$ 范畴。
\begin{proof}[命题~\ref*{st:10.13}的证明概要]
证明策略基本上与\stref{引理}{10.12}相同。一维模空间 $\mathcal H^T(y_0,\ldots,y_d)^{[0,1],1}$ 的边界点出现于以下情形：(i) 离根最近的顶点 $v$ 满足 $s_v=0$；(ii) 离叶子最近的某个顶点 $v$ 满足 $s_v=1$；或 (iii) 条件\eref{10.16}中的某个不等式成为等式。在前两种情形中，像\sr{fig:10.4}{图~10.4}那样将树分成若干部分，便得出相关代数贡献恰好是 $\Ai$ 函子方程\eref{1.6}中的两项。情形 (iii) 的贡献仍与 Gromov 紧化引入的额外点相抵消。
\SourceFigure{figure-10-4}{0.65215314}{10.4}
当然，这个形式上的纲要并不构成完整证明。不过，其中的大部分分析并不比定义 Fukaya 范畴时使用的分析复杂多少。因此，我们只讨论横截性问题的一个方面，即促使我们引入延迟函数的那个方面。为更清楚地说明这一点，暂时只考虑平凡延迟函数 $\sigma_k(r,s)=s$。

首先，可以一般地选择\eref{10.13}中的扰动数据，使所有空间 $\M^{d+1}(y_0,\ldots,y_d)^{[0,1]}$ 都正则。剩下的问题，是保证 $\mathcal H^T(y_0,\ldots,y_d)^{[0,1]}$ 的边界面仍然正则；这些边界面是乘积
\begin{equation}\label{eq:10.18}
\prod_{v\in\operatorname{Ve}(T)}\M^{|v|}(\widetilde y_{f_0(v)},\ldots,\widetilde y_{f_{|v|-1}(v)})^{[0,1]}
\end{equation}
中由等式
\begin{equation}\label{eq:10.19}
s_{v^+(e)}=s_{v^-(e)}
\end{equation}
截出的子集。只需要考虑维数不超过 $1$ 的模空间，使问题稍微容易一些。例如，在零维情形中，
\SourcePage{157}
所需陈述是：若 $\{u_v:S_v\to M\}\in\mathcal H^T(y_0,\ldots,y_d)^{[0,1]}$ 满足每个 $(r_v,s_v,u_v)$ 都是\eref{10.18}中相应因子的孤立点，则对任何边 $e$，\eref{10.19}都不成立。标准横截性论证说明，除非对应于相关顶点的两个映射相同，即
\begin{equation}\label{eq:10.20}
S_{v^-(e)}\cong S_{v^+(e)},\qquad u_{v^-(e)}=u_{v^+(e)},
\end{equation}
否则可以通过一般地选择扰动数据实现这一点。不幸的是，与前面对映射 $\Gamma$ 的讨论中相应的步骤不同，这个剩余情形根本无法排除，因为存在正则的伪全纯 $(d+1)$ 边形，其入向端的极限与一个或多个出向端的极限相同。现在允许非平凡的延迟函数；此时不再考虑\eref{10.19}，而是考虑\eref{10.16}中某些不等式成为等式的情形。一般地会有如下性质：\emph{若 $(r,s,u)$ 是某个模空间 $\M^{d+1}(y_0,\ldots,y_d)^{[0,1]}$（$d\geq2$）的孤立点，则延迟函数对所有 $k=1,\ldots,d$ 都满足 $\sigma_k(r,s)\ne s$。}这个条件显然排除了\eref{10.20}这样的病态行为，从而可以完成横截性论证。\znote{149}{病态情形\eref{10.20}的实质是：两个相邻顶点的参数相等时，它们的解可以拼成同一个圆盘，于是所谓“边界层”并不真的落在边界上，横截性论证落空。延迟函数 $\sigma_k$ 的作用正是把上、下圆盘的参数故意错开，使 $\sigma_k(r,s)\ne s$ 恒成立；这样\eref{10.19}的等号条件在任何孤立点上都不出现，横截性论证才得以完成。注意 $\sigma_k$ 只改变族沿 $s$ 方向的参数化，不改变任何模空间本身，故不影响已证得的结构。[来源：原书 §(10e) 关于\eref{10.13}、\eref{10.16}、\eref{10.20} 的讨论。]}对于一维模空间，平行的讨论只略复杂一些，这里略去。
\end{proof}
\begin{note}\label{st:10.14}
Floer 同调的同伦方法已经出现在 Floer 最初的论文\bcite{44}中。向 Fukaya 范畴的推广遵循 Fukaya–Oh–Ohta–Ono\bcite{60}的定理 15.19，不过我们以更明确的横截性论证替代了虚扰动理论。顺带指出，在非正合情形的 Floer 同调中，横截性问题更为严重，目前只找到间接的解决办法；见\bcite{94}，或参见\bcite{78}中作为模型的 Morse–Novikov 理论。
\end{note}
